\lstdefinelanguage{yaml}{
  morekeywords={true,false,null,y,yes,no,on,off},
  sensitive=false,
  morecomment=[l]{\#},
  morestring=[b]",
  morestring=[b]',
  alsoletter={-:},
  moredelim=[is][\color{black}]{\{}{\}},
  moredelim=[is][\color{black}]{[}{]},
  moredelim=[s][\itshape]{:}{\ },
  moredelim=[s][\itshape]{-}{\ }
}

\chapter{Sesión 1: Infraestructura base con Docker}
\label{ch:sesion1}

\section{Objetivo de la sesión}

Levantar un entorno de desarrollo reproducible con todos los servicios de infraestructura que NotiUAM necesita, orquestados con Docker Compose:
\begin{itemize}[leftmargin=*]
  \item \textbf{PostgreSQL 16} para persistencia relacional.
  \item \textbf{Redis 7} para caché y datos temporales.
  \item \textbf{Apache Kafka 3.7} (modo KRaft, sin Zookeeper) como broker de eventos.
  \item \textbf{Kafka UI} para inspeccionar topics y mensajes desde el navegador.
\end{itemize}

\section{Contexto}

El protocolo describe una arquitectura pub-sub con Kafka como broker central, PostgreSQL y Redis para persistencia y caché, y contenerización con Docker para reproducibilidad (secciones 2.3 y 6.2.3). Antes de escribir una sola línea de código Java o Dart, necesitábamos poder levantar esa infraestructura con un solo comando.

\section{Desarrollo}

\subsection{Estructura del repositorio}

Se creó el siguiente árbol en la raíz del proyecto:

\begin{lstlisting}[language=bash]
notiuam/
|- docker-compose.yml
|- .env
|- .gitignore
|- README.md
|- backend
\end{lstlisting}


\subsection{Variables de entorno (\texttt{.env})}

Se centralizaron credenciales y puertos en un archivo \texttt{.env} para no incrustarlos en el compose:

\begin{lstlisting}[language=bash]
POSTGRES_DB=notiuam
POSTGRES_USER=notiuam
POSTGRES_PASSWORD=notiuam_dev_pass
POSTGRES_PORT=5432

REDIS_PORT=6379

KAFKA_PORT=9092
KAFKA_UI_PORT=8080

BACKEND_PORT=8081
JWT_SECRET=change_me_in_prod_please_use_a_very_long_random_secret_at_least_256_bits
JWT_EXPIRATION_MS=3600000
JWT_REFRESH_EXPIRATION_MS=604800000
\end{lstlisting}

\begin{aviso}[Aviso]
El \texttt{JWT\_SECRET} de ejemplo \textbf{no debe usarse en producción}. Se reemplazará por un secreto de al menos 256 bits generado aleatoriamente.
\end{aviso}

\subsection{\texttt{docker-compose.yml}}

Se definieron cuatro servicios en una red interna \texttt{notiuam-net}, con volúmenes persistentes para cada uno.

\begin{lstlisting}[language=bash]
networks:
  notiuam-net:
    driver: bridge

volumes:
  postgres_data:
  redis_data:
  kafka_data:

services:

  postgres:
    image: postgres:16-alpine
    container_name: notiuam-postgres
    restart: unless-stopped
    environment:
      POSTGRES_DB: ${POSTGRES_DB}
      POSTGRES_USER: ${POSTGRES_USER}
      POSTGRES_PASSWORD: ${POSTGRES_PASSWORD}
    ports:
      - "${POSTGRES_PORT}:5432"
    volumes:
      - postgres_data:/var/lib/postgresql/data
    networks: [notiuam-net]
    healthcheck:
      test: ["CMD-SHELL", "pg_isready -U ${POSTGRES_USER} -d ${POSTGRES_DB}"]
      interval: 10s
      timeout: 5s
      retries: 5

  redis:
    image: redis:7-alpine
    container_name: notiuam-redis
    restart: unless-stopped
    command: ["redis-server", "--appendonly", "yes"]
    ports:
      - "${REDIS_PORT}:6379"
    volumes:
      - redis_data:/data
    networks: [notiuam-net]
    healthcheck:
      test: ["CMD", "redis-cli", "ping"]
      interval: 10s
      timeout: 5s
      retries: 5

  kafka:
    image: apache/kafka:3.7.0
    container_name: notiuam-kafka
    restart: unless-stopped
    ports:
      - "${KAFKA_PORT}:9092"
    environment:
      KAFKA_NODE_ID: "1"
      KAFKA_PROCESS_ROLES: "broker,controller"
      KAFKA_LISTENERS: "PLAINTEXT://0.0.0.0:9092,CONTROLLER://0.0.0.0:9093"
      KAFKA_ADVERTISED_LISTENERS: "PLAINTEXT://localhost:${KAFKA_PORT}"
      KAFKA_LISTENER_SECURITY_PROTOCOL_MAP: "CONTROLLER:PLAINTEXT,PLAINTEXT:PLAINTEXT"
      KAFKA_INTER_BROKER_LISTENER_NAME: "PLAINTEXT"
      KAFKA_CONTROLLER_LISTENER_NAMES: "CONTROLLER"
      KAFKA_CONTROLLER_QUORUM_VOTERS: "1@kafka:9093"
      KAFKA_OFFSETS_TOPIC_REPLICATION_FACTOR: "1"
      KAFKA_TRANSACTION_STATE_LOG_REPLICATION_FACTOR: "1"
      KAFKA_TRANSACTION_STATE_LOG_MIN_ISR: "1"
      KAFKA_GROUP_INITIAL_REBALANCE_DELAY_MS: "0"
      KAFKA_AUTO_CREATE_TOPICS_ENABLE: "true"
      KAFKA_NUM_PARTITIONS: "3"
    volumes:
      - kafka_data:/var/lib/kafka/data
    networks: [notiuam-net]
    healthcheck:
      test: ["CMD-SHELL", "/opt/kafka/bin/kafka-cluster.sh cluster-id --bootstrap-server localhost:9092 || exit 1"]
      interval: 15s
      timeout: 10s
      retries: 10

  kafka-ui:
    image: provectuslabs/kafka-ui:latest
    container_name: notiuam-kafka-ui
    restart: unless-stopped
    depends_on:
      kafka:
        condition: service_healthy
    ports:
      - "${KAFKA_UI_PORT}:8080"
    environment:
      KAFKA_CLUSTERS_0_NAME: notiuam-local
      KAFKA_CLUSTERS_0_BOOTSTRAPSERVERS: kafka:9092
      KAFKA_CLUSTERS_0_READONLY: "false"
    networks: [notiuam-net]
\end{lstlisting}

\subsection{Decisión técnica: imagen de Kafka}

\begin{aviso}[Cambio respecto al plan original]
Inicialmente se usó \texttt{bitnami/kafka:3.7}. Sin embargo, \textbf{Bitnami retiró su catálogo público de Docker Hub en septiembre de 2025}, por lo que el \texttt{docker pull} fallaba con \texttt{manifest unknown}. Se migró a la imagen oficial \texttt{apache/kafka:3.7.0}, mantenida por la Apache Software Foundation.

El cambio implicó ajustar los nombres de las variables de entorno:
\begin{itemize}[leftmargin=*]
  \item Bitnami usaba el prefijo \texttt{KAFKA\_CFG\_*} (p.\,ej. \texttt{KAFKA\_CFG\_NODE\_ID}).
  \item Apache usa el prefijo \texttt{KAFKA\_*} (p.\,ej. \texttt{KAFKA\_NODE\_ID}).
\end{itemize}
\end{aviso}

\subsection{Modo KRaft en lugar de Zookeeper}

Kafka 3.7 permite operar sin Zookeeper mediante \textbf{KRaft} (Kafka Raft). Se eligió este modo porque:
\begin{itemize}[leftmargin=*]
  \item Reduce el número de contenedores (uno en lugar de dos).
  \item Es el modo recomendado desde Kafka 3.0 y el único soportado a partir de Kafka 4.x.
  \item Simplifica el arranque en un solo nodo combinado \textit{broker + controller}.
\end{itemize}

\section{Verificación}

\subsection{Levantar los servicios}

\begin{lstlisting}[language=bash]
cd ~/Escritorio/notiuam
cp .env .env.example   # solo la primera vez, para tener un ejemplo
docker compose up -d
docker compose ps
\end{lstlisting}

\begin{exito}[Resultado esperado]
Cuatro contenedores en estado \texttt{Up (healthy)}:
\begin{lstlisting}[language=bash]
NAME                STATUS
notiuam-postgres    Up (healthy)
notiuam-redis       Up (healthy)
notiuam-kafka       Up (healthy)
notiuam-kafka-ui    Up
\end{lstlisting}
\end{exito}

\subsection{PostgreSQL}

\begin{lstlisting}[language=bash]
docker exec -it notiuam-postgres psql -U notiuam -d notiuam -c "\l"
\end{lstlisting}

Debe listar la base \texttt{notiuam}.

\subsection{Redis}

\begin{lstlisting}[language=bash]
docker exec -it notiuam-redis redis-cli ping
# Respuesta esperada: PONG
\end{lstlisting}

\subsection{Kafka}

\begin{lstlisting}[language=bash]
# Crear topic de prueba
docker exec -it notiuam-kafka /opt/kafka/bin/kafka-topics.sh \
  --bootstrap-server localhost:9092 \
  --create --topic test-topic --partitions 1 --replication-factor 1

# Listar topics
docker exec -it notiuam-kafka /opt/kafka/bin/kafka-topics.sh \
  --bootstrap-server localhost:9092 --list
# Debe aparecer: test-topic
\end{lstlisting}

\subsection{Kafka UI}

Abrir \url{http://localhost:8080} en el navegador. Debe aparecer el cluster \texttt{notiuam-local} y el topic \texttt{test-topic} creado en el paso anterior.

\section{Resultado}

\begin{exito}[Estado al cierre]
\begin{itemize}[leftmargin=*]
  \item Infraestructura base funcionando en Docker.
  \item Volúmenes persistentes para PostgreSQL, Redis y Kafka.
  \item Healthchecks configurados para detectar servicios listos.
  \item Archivo \texttt{.gitignore} y \texttt{README.md} en su lugar.
\end{itemize}
\end{exito}

\textbf{Pendiente para la siguiente sesión:} crear el proyecto backend con Spring Boot y conectarlo a estos servicios.